\section{Educational Demonstration: Making Sequential Decisions Visible}
\label{sec:educational-demo}

\subsection{Purpose and learning objectives}

To connect the preceding mathematics to observable decisions, we developed
an interactive game, \emph{You vs Bayesian Optimisation}. A human player
and a Bayesian optimisation (BO) algorithm search the same hidden
function, with the objective of obtaining the highest function value
within a fixed evaluation budget. The game is available at
\url{https://owenmarschner.github.io/you-vs-bo/}.

The intended audience is undergraduate mathematics students familiar with
basic probability, calculus and linear algebra, who have encountered
Gaussian processes and expected improvement through the preceding
sections. The demonstration has three learning objectives:
\begin{enumerate}
    \item interpret the distinction between a model's predictions,
    its uncertainty and the function values actually observed;
    \item explain how expected improvement converts predictions and
    uncertainty into a choice of the next input;
    \item recognise why a limited evaluation budget does not guarantee
    discovery of the global maximum.
\end{enumerate}

Playing the game makes the decision problem concrete. Before seeing
the hidden function, players must commit to an input using incomplete
information. The subsequent replay allows them to examine
BO's decisions under the same restriction. The emphasis is therefore on
explaining decisions using the information available when they were made,
rather than judging them solely with hindsight.


\subsection{Game structure and information available}
\label{subsec:demo-rules}

Both players begin with the same three observations, at
\[
    x=0.1,\qquad x=0.5,\qquad x=0.9.
\]
Each then receives five additional guesses by default, giving a total
budget of eight evaluations per player. The number of additional
guesses can be changed before a match, from three to ten.

On each turn, the human selects an input on the plot, where guesses are
spaced at intervals of \(0.01\). The \emph{Evaluate guess} button
evaluates the function at that input and displays the result. BO then
performs one evaluation using its own observation history. The players' subsequent observations
are kept separate: BO does not learn from the human's guesses, and the
human cannot see BO's chosen inputs, results or best value until the
match ends.

Hiding BO's inputs as well as its outputs is important. Even without
seeing the corresponding function values, the human could otherwise
infer which regions the algorithm considered promising and use that
information in later guesses. During play, the human's plot therefore
contains only their own observations. These points are not joined by
lines, since the function values between them remain unknown.

For player \(p\), let \(\mathcal{D}_p\) denote all observations available
to that player at the end of the match. Their score is
\[
    S_p=\max\{y:(x,y)\in\mathcal{D}_p\}.
\]
Thus the score includes the shared initial observations and depends on
the best value found at any stage, rather than the final guess or the
largest predicted value. The higher score wins; differences no greater
than \(10^{-9}\) are treated as a draw. Once both budgets are exhausted,
the function and both sets of observations are revealed.

This is a comparison between human intuition and the combination of a
GP model with an EI decision rule. The human is not supplied with a GP
during play, so the game does not isolate the effect of the acquisition
function alone.

\begin{figure}[htbp]
    \centering
    \IfFileExists{figures/demo-game.png}{
        \includegraphics[width=0.9\linewidth]{figures/demo-game.png}
    }{
        \fbox{\parbox[c][3.5cm][c]{0.85\linewidth}{
            \centering
            Game screenshot:
            upload \texttt{figures/demo-game.png}.
        }}
    }
    \caption{The game during the human player's search. The plot shows
    observed function values, while the score and remaining budget
    summarise progress. Selecting an input does not reveal its value
    until evaluation is confirmed. BO's search locations and results
    remain hidden.}
    \label{fig:demo-game}
\end{figure}


\subsection{Mathematical model and implementation}
\label{subsec:demo-numerics}

The demonstration considers continuous, noise-free functions on
\([0,1]\). Synthetic objectives are constructed from three smooth peaks:
\[
    f(x)=\sum_{j=1}^{3}
    a_j\exp\left(
        -\frac{(x-c_j)^2}{2w_j^2}
    \right),
\]
where \(a_j\), \(c_j\) and \(w_j\) specify each peak's amplitude, centre
and width. A seed determines these parameters, so entering the same
seed reproduces the same hidden function. Different seeds vary the peak
locations, heights and widths, although overlapping components can
produce fewer than three distinct local maxima. Before playing, the
human is told that the function is smooth, can have several peaks of
different widths and typically has high values in the hundreds.

These functions are inexpensive to evaluate. Their role is to simulate
the information restriction of an expensive experiment: the players
receive only a small number of evaluations, even though the software
could calculate other values. The function definition is not supplied
to BO's acquisition-selection routine.

BO uses the GP and noise-free EI framework described earlier
\cite{frazier2018tutorial}. Its prior has zero mean and a
squared-exponential kernel with length scale \(0.15\) and prior
standard deviation \(200\) at each input. The length scale implies a
prior correlation of about \(0.61\) between function values at inputs
that are \(0.15\) apart. These hyperparameters remain fixed throughout
a match. At each turn, the posterior is conditioned on BO's observations and EI is
calculated relative to BO's best observed value. The next input
maximises this acquisition score approximately.

Specifically, acquisition maximisation uses the candidate grid
\[
    \mathcal{G}
    =
    \left\{
        \frac{i}{2000}:i=0,\ldots,2000
    \right\}.
\]
Previously sampled BO inputs are excluded, and exact numerical ties
are resolved by selecting the lowest input. This produces a transparent,
reproducible approximation to continuous acquisition maximisation.
Calculating predictions and EI at the grid points does not consume
objective evaluations: only evaluating the selected input does.

GP conditioning is implemented through Cholesky solves, with diagonal
jitter equal to \(10^{-10}\) times the signal variance for numerical
stability. The model treats function evaluations as noise-free.
The implementation also handles the zero-uncertainty limit of EI and
clamps small negative variances caused by floating-point rounding.
A numerical failure preserves the previous match state rather than
silently replacing BO with another search strategy.

Checks of the mathematical implementation include comparison of GP
predictions with independently calculated small conditioning examples,
and comparison of EI with numerical integration. Additional tests check
equal budgets, independent observation histories, reproducibility and
scoring. These establish aspects of software correctness; they do not
by themselves establish educational effectiveness.


\subsection{Explaining BO through a decision replay}
\label{subsec:demo-decisions}

The competitive and explanatory stages are separated deliberately.
Showing BO's posterior and acquisition function during the match would
provide the human with information derived from BO's private
observations. Instead, the completed game provides a step-by-step
replay of the algorithm's reasoning.

Before each recorded BO evaluation, the replay displays the observations
available at that time, the GP posterior mean, and a shaded pointwise
\(95\%\) posterior band. Below this, the EI curve is shown with the
selected input marked. Numerical readouts give the predicted mean,
posterior standard deviation and EI at that input. This connects the
geometrical picture to the quantities used in the acquisition rule
\cite{frazier2018tutorial}.

The viewer can then reveal the evaluation. The new observation is added
and the posterior updates. The lower plot retains the original
pre-decision EI curve, making it possible to distinguish why the input
was selected from what was learned afterwards. Both plots use fixed
vertical scales across the replay so that viewers can compare turns
directly. The true function can be shown as a reference overlay and is
initially hidden within the replay.

The shaded band requires careful interpretation. It describes pointwise
posterior uncertainty under the chosen GP assumptions; it is neither a
guaranteed error bound nor a statement that the entire function lies
inside the band with probability \(0.95\).

The replay lets viewers compare the roles of predicted value and
uncertainty. An input with a mean below the best observed value can
still have high EI if its posterior assigns enough probability to
substantially better values. An input near a promising observation can
also have high EI when its predicted mean is high, even if its
uncertainty is lower.

\begin{figure}[htbp]
    \centering
    \IfFileExists{figures/demo-replay.png}{
        \includegraphics[width=0.95\linewidth]{figures/demo-replay.png}
    }{
        \fbox{\parbox[c][4.5cm][c]{0.88\linewidth}{
            \centering
            Replay screenshot:
            upload \texttt{figures/demo-replay.png}.
        }}
    }
    \caption{The BO decision replay before evaluation. The upper panel
    shows the posterior mean, observations and pointwise \(95\%\)
    posterior band. The lower panel shows EI and marks the selected
    input. Its outcome is revealed separately.}
    \label{fig:demo-replay}
\end{figure}


\subsection{Interpretation and limitations}
\label{subsec:demo-limitations}

The demonstration makes the consequences of sequential decisions
visible. Each new observation changes the information available to BO,
which can alter both its predictions and its next acquisition maximiser.
The replay connects these changes to the mathematical framework:
a selected input reflects the possibility of improvement under the
current posterior, while its eventual function value determines whether
that possibility is realised. A decision with high expected improvement
can therefore produce no improvement in the observed score.

Several limitations affect the interpretation of the game. Its small
evaluation budget can leave important regions unexplored, and its fixed
kernel hyperparameters determine which function shapes the model
considers plausible. In particular, inappropriate assumptions about
length scale can produce misleading predictions and uncertainty.
Furthermore, acquisition maximisation takes place on a finite grid,
so the selected input only approximates a continuous acquisition
maximiser. Revealing the hidden function helps identify missed features,
but neither the winning score nor the reveal certifies that a player
has found the global maximum.

The demonstration is consequently an explanatory example rather than
a performance benchmark. Assessing when BO is useful requires repeated
comparisons that control the initial observations and evaluation
budgets, account for variability, and examine sensitivity to modelling
choices. Section~5 develops this assessment by comparing EI, UCB and
uniform random search on contrasting test functions.
